\section*{Capítulo \ref{ch:g3:sound-around-us} --- O som ao nosso redor}

\begin{solution}{exo:g3:sound-around-us:1}
Um tremor de vaivém muito rápido. Violão: a corda tocada. Tambor: a
pele golpeada. Voz: as duas dobrinhas da garganta. Mosca: as asas
batendo.
\end{solution}

\begin{solution}{exo:g3:sound-around-us:2}
Que uma pele que soa é uma pele que treme: os grãos pulam porque a
pele debaixo deles sobe e desce depressa --- mesmo quando apenas um
grito, e não um toque, a põe em movimento.
\end{solution}

\begin{solution}{exo:g3:sound-around-us:3}
As mãos param o tremor do prato, e sem \omterm{def:g3:sound-around-us:vibration}{vibração} não há \omterm{def:g3:sound-around-us:sound}{som} --- é a
lei ``sem \omterm{def:g3:sound-around-us:vibration}{vibração}, sem \omterm{def:g3:sound-around-us:sound}{som}''.
\end{solution}

\begin{solution}{exo:g3:sound-around-us:4}
Puxe com mais \omterm{def:g2:pushes-and-pulls:force}{força}: o elástico balança mais largo --- um borrão mais
largo --- e o \omterm{def:g3:sound-around-us:sound}{som} fica mais forte. Tremor maior, \omterm{def:g3:sound-around-us:sound}{som} maior.
\end{solution}

\begin{solution}{exo:g3:sound-around-us:5}
A metade curta apertada; a flautinha; a corda mais fina. Coisas
curtas, pequenas e finas tremem mais depressa e cantam mais agudo.
\end{solution}

\begin{solution}{exo:g3:sound-around-us:6}
Eles sentem um zumbido --- as dobrinhas da garganta vibrando. O
zumbido para no instante em que o cantarolar para: sem \omterm{def:g3:sound-around-us:vibration}{vibração}, sem
\omterm{def:g3:sound-around-us:sound}{som}.
\end{solution}

\begin{solution}{exo:g3:sound-around-us:7}
As mãos batendo põem o \omterm{def:g2:air-around-us:air}{ar} a tremer; o tremor se espalha pelo \omterm{def:g2:air-around-us:air}{ar} como
os anéis num lago; ele chega ao ouvido e faz tremer a pelinha de lá
--- e a palma é ouvida.
\end{solution}

\begin{solution}{exo:g3:sound-around-us:8}
Sons muito altos sacodem a pelinha do ouvido --- o tambor de grãos de
arroz do próprio ouvido, mais fina que papel --- com força suficiente
para estragá-la para sempre; ela nunca cresce de novo.
\end{solution}

\begin{solution}{exo:g3:sound-around-us:9}
Os dentes curtos tocam as notas agudas e os dentes compridos tocam as
graves: coisas curtas tremem mais depressa e cantam mais agudo.
\end{solution}

\begin{solution}{exo:g3:sound-around-us:10}
Quem treme é o \omterm{def:g2:air-around-us:air}{ar} de dentro da garrafa --- o seu sopro o põe a
estremecer. Pôr mais água deixa uma coluna de \omterm{def:g2:air-around-us:air}{ar} \emph{mais curta}
acima da água, e curto canta agudo: é por isso que o \omterm{def:g3:sound-around-us:sound}{som} sobe. Se o
próprio vidro fosse o cantor, acrescentar água quase não mudaria a
nota.
\end{solution}

\begin{solution}{exo:g3:sound-around-us:11}
O \omterm{def:g3:sound-around-us:sound}{som} precisa de uma \emph{\omterm{def:g3:sound-around-us:vibration}{vibração}}: um tremor de vaivém rápido o
bastante para sacudir o \omterm{def:g2:air-around-us:air}{ar}, como a pele de tambor que fez o arroz
dançar. As imagens de uma tela no mudo mudam, mas nada ali faz o \omterm{def:g2:air-around-us:air}{ar}
tremer --- cores que mudam não sacodem nada, então não há nada para
um ouvido pegar.
\end{solution}
